% Source: 04 -Mon Oct 5 2026.pdf, page 5. Content remains in source order.
\sourcepage{04}{5}
\sourcefigure[0.85]{assets/04/page-005-circuit-configuration.png}
The circuit implements a boolean function $y=C(x_1,x_2,\ldots,x_n)$.
Y (the output) is the result of the values given in input to the function ( the circuit) 

For a circuit with $n$ inputs, a truth assignment is a vector with one bit per input $\vec b\in\bits^n$.

\subsection{Configuration} 

After the circuit has commuted every wire in the circuit will carry a stable boolean value.

The set of boolean values carried by all wires of a circuit $C$ under a truth-assignment $\vec b$ is called Configuration $\gamma_C$ of $C$ under $\vec b$. 

Since the input is part of the configuration, and the truth assignemnt has N bits, there are 2 to the N possible truth assignments and therefore 2 to the N possible configurations. 

\textbf{DEF:} A boolean circuit $C(x_1,x_2,\ldots,x_n)$ is satisfiable if there exists a truth assignment such that the output of the circuit has value 1 :  $\vec b\in\bits^n:\ C(\vec b)=1$.
